%% A281_Traeger_und_Information_En.tex
%% FFGFT A-Series -- A281: Carrier and Information
%% Autor: Johann Pascher  --  ORCID 0009-0000-6518-4064
%% Language: English  --  Engine: LuaLaTeX
%% Source: landauer_erweitert.tex (undated)
%% Companion document to A280

\documentclass[11pt,a4paper]{article}

% ---- Engine: LuaLaTeX ----
\usepackage{fontspec}
\usepackage{mathtools}
\usepackage{unicode-math}
\usepackage{polyglossia}
\setmainlanguage{english}

\usepackage{microtype}
\emergencystretch=3.5em
\tolerance=2500
\hbadness=3000
\usepackage{amsthm}
\usepackage{enumitem}
\usepackage{booktabs}
\usepackage{longtable}
\usepackage{array}
\usepackage[a4paper,margin=2.6cm,headsep=1.1em,headheight=22pt]{geometry}
\usepackage{fancyhdr}
\usepackage{titlesec}
\usepackage{xcolor}
\usepackage{needspace}
\usepackage{tocloft}

% ---- Fonts ----
\setmainfont{Inter}[
  Scale=1.02,
  Ligatures=TeX ]
\setsansfont{Inter}[Scale=MatchLowercase,Ligatures=TeX]
\setmonofont{JetBrains Mono}[Scale=0.88]
\setmathfont{Libertinus Math}

% ---- Colours ----
\definecolor{accent}{HTML}{1F3A5F}
\definecolor{accent2}{HTML}{6A4E2E}
\definecolor{rulecol}{HTML}{B9C2CE}
\definecolor{linknav}{HTML}{2C5F8A}
\definecolor{linkcite}{HTML}{7A5C2E}

% ---- Schichtmarkierungen (A-Serie) ----
\definecolor{laySet}{HTML}{7A2E2E}
\definecolor{layK}{HTML}{1D5C3A}
\definecolor{layB}{HTML}{1F3A5F}
\definecolor{layS}{HTML}{6A4E2E}
\definecolor{layH}{HTML}{5A2E6A}

\newcommand{\LS}{\textcolor{laySet}{\textbf{\footnotesize\textsf{[STIPULATION]}}}\,}
\newcommand{\LK}{\textcolor{layK}{\textbf{\footnotesize\textsf{[K]}}}\,}
\newcommand{\LB}{\textcolor{layB}{\textbf{\footnotesize\textsf{[B]}}}\,}
\newcommand{\LSk}{\textcolor{layS}{\textbf{\footnotesize\textsf{[S]}}}\,}
\newcommand{\LH}{\textcolor{layH}{\textbf{\footnotesize\textsf{[H]}}}\,}

% ---- Theorem-Umgebungen ----
\newtheoremstyle{ffgft}{0.8em}{0.4em}{\itshape}{}{\bfseries\color{accent}}{.}{0.5em}{}
\newcommand{\qedsq}{\raisebox{0.03em}{\framebox[0.62em]{\rule{0pt}{0.42em}}}}
\theoremstyle{ffgft}
\newtheorem{definition}{Definition}
\newtheorem{satz}{Proposition}

% ---- Hyperref ----
\usepackage[
  colorlinks=true,
  linkcolor=linknav,
  citecolor=linkcite,
  urlcolor=linknav,
  bookmarksnumbered=true,
  bookmarksopen=true,
  pdfstartview=FitH ]{hyperref}
\hypersetup{
  pdftitle={Carrier and Information (A281)},
  pdfauthor={Johann Pascher},
  pdfsubject={Landauer's principle --- textual criticism and reception history},
  pdfkeywords={Landauer, carrier, information, hypostatisation, reception, FFGFT, A281}
}

% ---- Section styling ----
\titleformat{\section}
  {\color{accent}\normalfont\Large\bfseries}
  {\color{accent2}\thesection}{0.7em}{}
\titleformat{\subsection}
  {\color{accent}\normalfont\large\bfseries}
  {\color{accent2}\thesubsection}{0.6em}{}
\titlespacing*{\section}{0pt}{1.6\baselineskip}{0.7\baselineskip}
\titlespacing*{\subsection}{0pt}{1.1\baselineskip}{0.4\baselineskip}

% ---- Headers ----
\pagestyle{fancy}
\fancyhf{}
\renewcommand{\headrulewidth}{0.4pt}
\renewcommand{\headrule}{\hbox to\headwidth{\color{rulecol}\leaders\hrule height \headrulewidth\hfill}}
\fancyhead[L]{\footnotesize\color{accent}\itshape FFGFT $\cdot$ A281 $\cdot$ Carrier and Information}
\fancyhead[R]{\footnotesize\color{accent2}\thepage}

% ---- TOC look ----
\renewcommand{\cftsecfont}{\normalfont}
\renewcommand{\cftsecpagefont}{\normalfont}
\renewcommand{\cftsecleader}{\cftdotfill{\cftdotsep}}
\setlength{\cftbeforesecskip}{3pt}

\setlength{\parskip}{0.35em}
\setlength{\parindent}{0pt}
\linespread{1.06}

\begin{document}

% ===================== TITLE =====================
\begin{titlepage}
\centering
\vspace*{2.2cm}
{\color{rulecol}\rule{0.5\textwidth}{0.6pt}}\\[1.4cm]
{\color{accent}\Huge\bfseries Carrier and Information\par}
\vspace{1.0cm}
{\large\itshape What Landauer proved in 1961\\ --- and what the reception made of it\par}
\vspace{0.6cm}
{\color{accent2}\rule{0.28\textwidth}{0.4pt}}\\[1.0cm]
{\large FFGFT A-Series $\cdot$ A281\par}
\vspace{0.4cm}
{\normalsize Companion document to A280 --- textual criticism, reception history\\ and a critique of language\par}
\vfill
{\large Johann Pascher\par}
\vspace{0.3cm}
{\footnotesize ORCID: 0009-0000-6518-4064\par}
\vspace{0.2cm}
{\footnotesize July 2026\par}
\vspace{1.4cm}
{\color{rulecol}\rule{0.5\textwidth}{0.6pt}}
\vspace*{1.5cm}
\end{titlepage}

% ===================== TOC =====================
\tableofcontents
\thispagestyle{fancy}
\clearpage

% ===================== ABSTRACT =====================
\begingroup\small\noindent
\textbf{Abstract.} The familiar claim that “erasing a bit produces at least
$kT\ln 2$ of heat” goes beyond what Landauer's 1961 paper proves. Landauer shows
something precise: if a physical two-state system in thermal equilibrium is forced
into a definite state irrespective of its initial state, at least $kT\ln 2$ must
be dissipated. That is a statement about \emph{carrier operations}. The step from
there to a universal claim about \emph{information} as such is not taken in
Landauer's paper; it is encouraged by a conceptual ambiguity in the word “bit”,
but it is not a proved consequence of his arguments. This paper separates the two,
documents Landauer's own restriction, traces the generalisation through the
reception literature, and names the linguistic mechanism --- the hypostatisation
of an abstraction --- that carried it. \emph{Relation to A280:} A281 adds nothing
to the physics. A280 does the accounting; A281 does the textual criticism. The
dictionary mapping the two vocabularies is given in
Section~\ref{sec:dictionary}.
\par\endgroup

\bigskip

\begingroup\small\noindent
\textbf{Layer markings.} \LS declared starting point (German corpus token:
[SETZUNG]) --- \LB derived algebraically/mathematically from declared foundations
--- \LK empirically tested or documented from primary sources --- \LSk
plausibility sketch --- \LH open research question.
\par\endgroup

\bigskip

% ===================== 1 =====================
\section{The question at issue}

Since Landauer's 1961 paper [B1] it has often been claimed in the literature that
every erased bit must produce at least $kT\ln 2$ of heat. The claim has become a
fixture of computer science, physics and the philosophy of science.

The present paper holds that this generalisation goes beyond what Landauer's proof
shows. Landauer makes a precise statement about physical carrier operations. The
jump from there to a universal statement about information as an abstract quantity
is not his step --- it was taken in the reception, encouraged by a conceptual
imprecision that is built into the discipline.

Two delimitations first, so that the reach of this paper is not overestimated.

\LB \textbf{First: this is not a refutation of Landauer's principle.} What is
contested here is the \emph{popular short form}, not Landauer's calculation. For
the physical operation Landauer considers, his result is correct. What is
contested is the jurisdiction of that result over a subject matter --- abstract
information --- about which it makes no statement.

\LB \textbf{Second: the physical substance is in A280, not here.} That the
thermodynamic lower bound attaches to the physical region and not to its
description is derived there from Liouville. A281 adds nothing to that. What A281
does is different and stands on its own: the textual analysis of Landauer 1961,
the documentation of Landauer's own restriction, the reception history of the
generalisation, and the critique of the language of digital technology.

\subsection{Dictionary: A280 and A281}
\label{sec:dictionary}

\LS The two documents draw the same distinction in different words. So that the
corpus does not end up carrying two parallel terminologies for one matter, the
mapping is fixed here and is binding:

\begin{center}
\small
\begin{tabular}{@{}p{0.36\textwidth}p{0.36\textwidth}p{0.20\textwidth}@{}}
\toprule
\textbf{A280 (accounting)} & \textbf{A281 (textual criticism)} & \textbf{Level} \\
\midrule
region in phase space & carrier & physical \\
description & information & abstract \\
region reduction $\{A,B\}\to Z$ & carrier operation \textup{RESTORE TO ONE} & physical \\
non-bijectivity of the region map & shrinking of the carrier state space & physical \\
descriptive decision & interpretive erasure & abstract \\
\bottomrule
\end{tabular}
\end{center}

\LB The mapping is a translation, not an extension. Whoever reads “carrier” may
think “region”, and conversely. A statement that holds in one column holds in the
other --- otherwise there is an error in one of the two documents, not a new
matter of fact.

% ===================== 2 =====================
\section{The core distinction}

\subsection{Definitions}

\begin{definition}[Information]
Information is an abstract quantity. A bit is a unit of information representing
one of two states ($0$ or $1$). Information has no mass, no energy, no entropy and
no physical location.
\end{definition}

\begin{definition}[Carrier]
The carrier of the information is a physical system that represents the
information: an electrical circuit, a magnetic spin, a photon or any other
physical two-state system. The carrier has energy, entropy and further physical
properties.
\end{definition}

\LS Both definitions are stipulations. They are not proved; they fix what is being
talked about in what follows. The entire yield of the paper depends on keeping the
separation intact.

\subsection{Energy attaches to the carrier}

\begin{satz}
No energy can be assigned to a bit of information, only to its carrier.
\end{satz}

\LB \textbf{Proof.} Suppose a bit of information had an intrinsic energy. Then that
energy would have to be independent of the physical realisation. But a bit with
the value $1$ can be represented by different physical systems with different
energies: by 5\,V in a CMOS circuit, by 1\,V in another technology, by a magnetic
spin with a given exchange energy, by a photon of a given frequency. Since the
information content is identical in all cases while the energy varies, the energy
cannot be a property of the information. It is a property of the carrier.
\hfill\qedsq

The argument is a multiple-realisability argument and carries exactly as far as
multiple realisability reaches --- that is, all the way: there is no class of
realisations in which the energy would co-vary with the information content.

\subsection{The bus example}

\LB An instructive example is the bit width of a data bus:

\begin{itemize}[leftmargin=1.4em,itemsep=0.25em]
  \item The \textbf{bus width} (say 64 bits) is a fixed physical property of the
        hardware. It determines how many carriers are switched simultaneously.
  \item The \textbf{information} is what is encoded by the patterns. The same bit
        pattern can mean different things in different contexts.
\end{itemize}

A 64-bit bus always switches 64 carriers, regardless of whether the information is
a number, a letter or a picture element. The lines have a certain capacitance and
require energy for recharging. The information does not have that energy.

% ===================== 3 =====================
\section{Landauer's model --- what it shows}

\subsection{The particle in the double well}

\LK Landauer's central model (Fig.~1 in [B1]) is a classical particle in a
bistable potential well:

\begin{itemize}[leftmargin=1.4em,itemsep=0.2em]
  \item left minimum: state $0$;
  \item right minimum: state $1$;
  \item the operation \textup{RESTORE TO ONE} forces the particle into the right
        well irrespective of its initial state.
\end{itemize}

This is a classical macroscopic model. It does not capture the full breadth of
modern quantum-information concepts, but it remains valid as an analogy for
bistable systems --- spin-½ systems are formally equivalent. What matters is what
the model does \emph{not} do: it shows under which physical conditions dissipation
is unavoidable. It does \emph{not} show that those conditions are met whenever a
bit is erased in the information-theoretic sense.

\begin{center}
\small
\begin{tabular}{@{}p{0.45\textwidth}p{0.45\textwidth}@{}}
\toprule
\textbf{What Landauer's model does} & \textbf{What Landauer's model does not do} \\
\midrule
proof for physical carrier operations & proof for information-theoretic erasure \\
lower bound for \textup{RESTORE TO ONE} & statement about abstract bits \\
validity for a thermalised ensemble & validity for arbitrary computational processes \\
physical entropy $S = k\ln W$ & bridge to Shannon entropy \\
\bottomrule
\end{tabular}
\end{center}

\subsection{The entropy calculation}

\LK Landauer computes the entropy change on resetting as
\begin{equation}
  \Delta S \;=\; k \ln 2 - k \ln 1 \;=\; k \ln 2 .
  \label{eq:landauer}
\end{equation}

This calculation presupposes:
\begin{enumerate}[leftmargin=1.6em,itemsep=0.2em]
  \item \textbf{before:} the carrier can take 2 physical states ($W=2$);
  \item \textbf{after:} the carrier can take only 1 physical state ($W=1$).
\end{enumerate}

Both presuppositions are statements about the carrier. Neither is a statement
about the information content.

% ===================== 4 =====================
\section{The objection and its reach}

\LB The entropy change in equation~\eqref{eq:landauer} is a property of the
\textbf{physical carrier}, not of the information. If a bit is erased
\emph{interpretively} --- by redefining its meaning, declaring the value invalid,
striking the address from a table --- then the number of physical states of the
carrier does \textbf{not} change. It still holds that $W = 2$, hence
\begin{equation}
  \Delta S \;=\; k \ln 2 - k \ln 2 \;=\; 0 .
  \label{eq:interpretive}
\end{equation}

\begin{satz}
Landauer's principle bites only if the operation physically shrinks the state
space of the carrier. A purely interpretive erasure does not shrink it.
\end{satz}

\LB \textbf{Proof.} Let $W$ be the number of physical microstates of the carrier.
An interpretive erasure declares one state no longer valid; it is an operation at
the level of description and leaves the physical state space unchanged. The
entropy of the carrier remains $S = k\ln W$. Only when the carrier itself is
forced into a single state ($W' = 1$) does the entropy change to $S' = 0$. That is
a physical operation on the carrier, not on the information content.
\hfill\qedsq

\subsection{What this objection is not directed against}

\LB Precision is needed here, otherwise the objection is aimed at an opponent who
does not exist. Equation~\eqref{eq:interpretive} does \emph{not} refute Landauer.
Landauer nowhere claims that a conceptual redesignation costs heat; nobody in the
technical literature claims that. What
equation~\eqref{eq:interpretive} shows is solely that the two processes ---
interpretive erasure and carrier reset --- come apart, and that only the second
falls under Landauer's bound.

That is not a small thing, but it is also not what it looks like at first glance.
The yield is a \emph{finding about jurisdiction}, not a refutation:

\begin{quote}
\LB Landauer's bound has jurisdiction over carrier operations that shrink the
state space. For everything else that ordinary language calls “erasing”, it has no
jurisdiction --- neither confirming nor refuting.
\end{quote}

\LB The real question thereby shifts, and it is harder than the original one:
\emph{does information-theoretic erasure enforce a state-space-shrinking carrier
operation?} Bennett [B2][B8] showed that in general it does not --- any logically
irreversible function can be executed reversibly if the intermediate states are
carried along. Only the final relinquishing of carrier states costs. So the answer
for the general case is negative, and it is an answer Landauer could not yet give
in 1961.

% ===================== 5 =====================
\section{Landauer's own restriction}

\LK Landauer formulated a restriction that the later reception mostly passes over:

\begin{quote}
\itshape
“Our reset operation, in the preceding discussion, was applied to a thermal
equilibrium ensemble. In actuality we would like to know what happens in a
particular computing circuit which will work on information which has not yet been
thermalized \ldots” \upshape [B1]
\end{quote}

\LK Landauer resolves the question himself immediately afterwards: for a random
sequence of ones and zeros he equates the statistical ensemble with the time
average and thereby arrives at the same dissipation per reset as for the
thermalised ensemble. For non-random data --- if, say, all inputs are already
\textup{ONE} --- he expressly concedes that no entropy change occurs.

Three things follow:

\begin{enumerate}[leftmargin=1.6em,itemsep=0.3em]
  \item \LK Landauer is aware of the limits of his model and works them out
        himself. The charge of over-generalisation hits the reception, not him.
  \item \LB His result $kT\ln 2$ holds strictly for the equilibrium ensemble. For
        correlated, non-random data, $kT\ln 2$ per carrier is an \emph{upper}
        bound; the applicable bound is $kT\,H$, with $H$ the entropy of the source
        in nats. This is not a weakening of the principle but its correct
        statement: the bound follows the actual state count, not the bit width.
  \item \LB Above all: this extension too remains a statement about
        \emph{carrier states}. The transition to abstract information is not made
        here either.
\end{enumerate}

% ===================== 6 =====================
\section{What Landauer proves --- and what he does not}

The popular short form reads:

\begin{quote}
\emph{“Every erased bit produces $kT\ln 2$ of heat.”}
\end{quote}

\LB This claim goes beyond Landauer's proof. What is proved is:

\begin{quote}
\emph{If a physical two-state system in thermal equilibrium is forced into a
definite state irrespective of its initial state, at least $kT\ln 2$ of heat must
be released. This dissipation is a property of the physical carrier operation.}
\end{quote}

\LK The step Landauer's proof does not contain is the bridge from the carrier
operation to abstract information. Landauer uses “bit” throughout in a double
sense --- sometimes for the physical object, sometimes for the unit of
information. This ambiguity is present in the original paper and encouraged the
later reception, in which a statement about carriers became a statement about
information.

\LB That every information-theoretic erasure requires a state-space-shrinking
carrier operation is presupposed, not derived --- and by Bennett it is in general
false.

% ===================== 7 =====================
\section{The modern thermodynamics of information}

\LK The thermodynamics of information with feedback [B5][B6] gives, for the
erasure process, the generalised bound
\begin{equation}
  Q_{\min} \;=\; k T\,(\ln 2 - I), \qquad I \in [0,\ \ln 2],
  \label{eq:feedback}
\end{equation}
where $I$ is the mutual information between measuring device and system state,
measured in \textbf{nats} and therefore dimensionless. Equivalently in the form of
Sagawa and Ueda: $W_{\text{ext}} \le -\Delta F + kT\,I$.

\begin{itemize}[leftmargin=1.4em,itemsep=0.25em]
  \item \LB If the experimenter knows whether the carrier is at $0$ or $1$
        ($I = \ln 2$), the carrier can be reset reversibly --- $Q_{\min} = 0$.
  \item \LB If not ($I = 0$), the reset must be irreversible ---
        $Q_{\min} = kT\ln 2$.
\end{itemize}

\LB The heat therefore arises not \emph{because} information is lost, but because
of the lack of knowledge about the carrier state. The information is the label
stuck onto the carrier; the price attaches to how far the label actually pins the
carrier down.

\LB And the symmetry must be seen at the same time: the knowledge $I$ is itself
stored on a carrier. Whoever uses it to reset reversibly has not avoided the price
but displaced it --- to the later moment at which the carrier of the knowledge is
itself reset. That is Bennett's resolution of Maxwell's demon [B8], and it
prevents equation~\eqref{eq:feedback} from being read as a circumvention of the
bound.

% ===================== 8 =====================
\section{The perpetuation in the reception}

\LK The conceptual imprecision would have remained inconsequential had the
subsequent literature corrected it. The opposite happened. The majority of
publications after 1961 adopt the principle in generalised form without going back
to Landauer's own restrictions.

\begin{itemize}[leftmargin=1.4em,itemsep=0.3em]
  \item \LK \textbf{Lairez} [B9] notes that criticism of the generalisation has
        been available for decades but has had no broad effect.
  \item \LK \textbf{Norton} [B4] shows that the standard argument rests on an
        untenable identification of dynamic and static probabilities;
        \textbf{Earman and Norton} [B3] had already worked through the line from
        Szilárd to Landauer in 1999.
  \item \LK \textbf{Hemmo and Shenker} [B10] show that in the light of the
        multiple-computations theorem the principle cannot be universally valid.
\end{itemize}

\LSk The reason for the persistence lies not in a refutation of this criticism but
in structural features of scientific communication: the formula $kT\ln 2$ is
catchy and easy to cite. The 1961 original is rarely read; what gets cited are
secondary sources citing secondary sources. Trust in the authority of a
distinguished industrial researcher does the rest. The consequence is that a step
never taken in the proof is traded as a “fundamental law of nature” --- a
description that does not appear in the original paper.

\LB This diagnosis is a plausibility sketch and is marked as such. It names
mechanisms that are individually observable but whose combined effect is not
quantified. A defensible version would be a citation analysis: the share of papers
citing [B1] that also state the ensemble restriction. That analysis is not
provided here.

% ===================== 9 =====================
\section{The imprecision in the language of digital technology}

\LB The problem is not confined to Landauer. The whole of digital technology uses a
language that systematically states carrier properties as properties of
information. The practice is so deeply entrenched that practitioners rarely
perceive it as problematic.

\begin{center}
\small
\begin{longtable}{@{}p{0.26\textwidth}p{0.34\textwidth}p{0.30\textwidth}@{}}
\toprule
\textbf{Customary expression} & \textbf{Physical reality (carrier)} & \textbf{Problematic implication} \\
\midrule
\endfirsthead
\toprule
\textbf{Customary expression} & \textbf{Physical reality (carrier)} & \textbf{Problematic implication} \\
\midrule
\endhead
“set a bit” & apply a voltage to a capacitor & information has a state \\
“erase a bit” & set the carrier to a reference voltage & information can be annihilated \\
“transmit a bit” & reconstruct a signal on a new line & information is transportable \\
“copy a bit” & read a voltage, generate it elsewhere & information multiplies itself \\
“the RAM holds 8\,GB” & 8 billion physical memory cells & information is a bulk quantity \\
“bit rate” & voltage changes per second & amount of information per unit time \\
\bottomrule
\end{longtable}
\end{center}

In each case a property of the carrier is phrased as if it were a property of the
information. In practice engineers know that they are handling voltages, currents
and charges. In the descriptive language, however, the bit becomes the acting
subject: it “flows through the bus”, “is stored”, “is erased”. This language is
convenient for many purposes --- but it tempts one to apply physical laws directly
to the abstraction, as the popular Landauer reception did.

\LB A more precise usage would distinguish:

\begin{itemize}[leftmargin=1.4em,itemsep=0.25em]
  \item \textbf{Not:} “the bit has energy.” \quad
        \textbf{But:} “the carrier of this bit has energy.”
  \item \textbf{Not:} “the bit is erased.” \quad
        \textbf{But:} “the carrier is set to a definite state.”
  \item \textbf{Not:} “the bit costs $kT\ln 2$.” \quad
        \textbf{But:} “this carrier operation costs at least $kT\ln 2$ under these
        conditions.”
\end{itemize}

The connection to Landauer is direct: he uses “bit” in exactly this double sense.
The ambiguity is not his invention; it is built into the discipline and it carried
the generalisation.

% ===================== 10 =====================
\section{Epistemological background: hypostatisation}

\LSk The mechanism described --- an abstraction is treated as a physical quantity
and equipped with physical properties --- has a name: \emph{hypostatisation} (or
reification), the transformation of an abstract concept into a supposedly real,
substantial object.

\LK The bit is an abstraction. Shannon [B11] defined it as a statistical property
of signals, not as a physical object. Wiener stated the same point explicitly:
information is information, not matter or energy [B12]. That clarity was lost in
engineering practice. The bit became an object to which energy, entropy and
location could be ascribed.

\LSk The mechanism is not new in the history of science. Entropy, the electric
field, the gene --- all of these abstractions were at certain times treated as
substantial things, with consequences for theory. In the case of information the
hypostatisation leads to a false physics: if the bit itself had entropy, then
erasing it would have to produce heat --- independently of the carrier. That
conclusion is compelling only for someone who has confused the abstraction with
physical reality.

\LB This section is marked as a sketch because it offers an interpretation and
provides no demonstration. The historical parallels are analogies; they support
the thesis, they do not prove it.

% ===================== 11 =====================
\section{Conclusion}

\begin{enumerate}[leftmargin=1.6em,itemsep=0.35em]
  \item \LB Landauer's proof for the physical operation \textup{RESTORE TO ONE} on
        a thermal equilibrium ensemble is correct and precise.
  \item \LB It shows a lower bound for \emph{carrier operations} --- not a
        universal statement about abstract information.
  \item \LB The transfer to the information-theoretic notion of the bit
        presupposes that every information-theoretic erasure is bound up with a
        state-space-shrinking carrier operation. Landauer does not prove this
        connection --- and by Bennett it does not hold in general.
  \item \LK The double sense of “bit” is present in the original paper and
        encouraged the generalisation in the reception.
  \item \LK Landauer himself discusses the limits of his ensemble argument and
        concedes that non-random input data lower the dissipation below
        $kT\ln 2$.
  \item \LB What is not shown here: that Landauer's principle is false. What is
        shown is that it has jurisdiction over a subject matter other than the one
        it is usually cited for.
\end{enumerate}

% ===================== 12 =====================
\section{Outlook}

\LB A formulation that reflects the actual proof reads:

\begin{quote}
\emph{The physical operation that puts a bistable system in thermal equilibrium
into a definite state irrespective of its initial state requires at least
$kT\ln 2$ of dissipation. If this system serves as an information carrier and the
operation corresponds to erasing a bit, the bound applies to the associated
carrier operation --- not to information-theoretic erasure as such.}
\end{quote}

\LH \textbf{Open question.} Under what conditions does a given computational
process enforce a state-space-shrinking carrier operation? Bennett answers the
limiting case (in principle never, if everything is carried along); the
practically interesting question is quantitative: how much carrier state space
must an architecture finally give up at a given memory depth? This is the same
open question as the continuous entropy current in A280 --- posed from the other
side.

% ===================== 13 =====================
\section{Layer-status summary}

\begin{center}
\small
\begin{longtable}{@{}p{0.72\textwidth}l@{}}
\toprule
\textbf{Statement} & \textbf{Status} \\
\midrule
\endfirsthead
\toprule
\textbf{Statement} & \textbf{Status} \\
\midrule
\endhead
Definition of information (abstract, without energy) & \LS \\
Definition of carrier (physical, with energy) & \LS \\
Dictionary mapping A280 $\leftrightarrow$ A281 & \LS \\
Proposition 1: energy attaches to the carrier, not to the bit & \LB \\
Bus example: bus width is hardware, not information & \LB \\
Landauer's double-well model and \textup{RESTORE TO ONE} & \LK \\
$\Delta S = k\ln 2$ presupposes $W{=}2 \to W{=}1$ at the carrier & \LK \\
Proposition 2: interpretive erasure does not shrink $W$ & \LB \\
The objection is a finding about jurisdiction, not a refutation & \LB \\
Bennett: logically irreversible functions executable reversibly & \LB \\
Landauer's ensemble restriction (original quotation) & \LK \\
Landauer's own resolution for random data & \LK \\
For correlated data the bound is $kT\,H$, not $kT\ln 2$ & \LB \\
Landauer uses “bit” in a double sense & \LK \\
Feedback bound $Q_{\min} = kT(\ln 2 - I)$, $I$ in nats & \LK \\
The knowledge $I$ is itself carrier-bound (Bennett) & \LB \\
Reception findings: Lairez / Norton / Hemmo--Shenker & \LK \\
Structural reasons for the persistence & \LSk \\
Language table: carrier versus information & \LB \\
Hypostatisation as an interpretive frame & \LSk \\
Historical parallels (entropy, field, gene) & \LSk \\
No demonstration that Landauer's principle is false & \LB \\
Quantitative form of the unavoidable state-space relinquishment & \LH \\
\bottomrule
\end{longtable}
\end{center}

\textbf{Balance:} 10$\times$ \LB $\cdot$ 8$\times$ \LK $\cdot$ 3$\times$ \LS
$\cdot$ 3$\times$ \LSk $\cdot$ 1$\times$ \LH $=$ 25 entries.

% ===================== APPENDIX =====================
% ===================== SOURCES =====================
\clearpage
\phantomsection
\addcontentsline{toc}{section}{Sources}
\section*{Sources}

\begingroup
\small
\begin{description}[leftmargin=3.2em,style=sameline,itemsep=0.3em]

\item[{[B1]}] R. Landauer, \emph{Irreversibility and Heat Generation in the
Computing Process}, IBM J.\ Res.\ Dev.\ \textbf{5}, 183 (1961).

\item[{[B2]}] C.~H. Bennett, \emph{Logical Reversibility of Computation}, IBM J.\
Res.\ Dev.\ \textbf{17}, 525 (1973).

\item[{[B3]}] J. Earman and J.~D. Norton, \emph{Exorcist XIV: The Wrath of
Maxwell's Demon. Part II: From Szilard to Landauer and Beyond}, Stud.\ Hist.\
Phil.\ Mod.\ Phys.\ \textbf{30}, 1 (1999).

\item[{[B4]}] J.~D. Norton, \emph{All Shook Up: Fluctuations, Maxwell's Demon and
the Thermodynamics of Computation}, Entropy \textbf{15}, 4432 (2013).

\item[{[B5]}] T. Sagawa and M. Ueda, \emph{Nonequilibrium Thermodynamics of
Feedback Control}, Phys.\ Rev.\ E \textbf{85}, 021104 (2012).

\item[{[B6]}] J.~M.~R. Parrondo, J.~M. Horowitz and T. Sagawa,
\emph{Thermodynamics of Information}, Nature Physics \textbf{11}, 131 (2015).

\item[{[B7]}] J. Bub, \emph{Maxwell's Demon and the Thermodynamics of
Computation}, Stud.\ Hist.\ Phil.\ Mod.\ Phys.\ \textbf{32}, 569 (2001).

\item[{[B8]}] C.~H. Bennett, \emph{The Thermodynamics of Computation --- a
Review}, Int.\ J.\ Theor.\ Phys.\ \textbf{21}, 905 (1982).

\item[{[B9]}] \emph{(incomplete --- identifier to be supplied)} D. Lairez,
\emph{Landauer's Principle: A Critical Assessment}, arXiv preprint (2024).

\item[{[B10]}] M. Hemmo and O. Shenker, \emph{The Multiple-Computations Theorem
and the Physics of Singling Out a Computation}, The Monist \textbf{105}, 175
(2022).

\item[{[B11]}] C.~E. Shannon, \emph{A Mathematical Theory of Communication}, Bell
System Technical Journal \textbf{27}, 379 (1948).

\item[{[B12]}] \emph{(reconstructed --- page number to be checked)} N. Wiener,
\emph{Cybernetics: or Control and Communication in the Animal and the Machine},
MIT Press, Cambridge, MA (1948).

\end{description}
\endgroup

\vspace{1em}
\noindent
{\small\itshape Entries marked \emph{incomplete} or \emph{reconstructed} were
supplied from the citation context or carry an open placeholder; they are to be
checked against the working files before publication.}

\end{document}
